% !PRINTED 20
\weberpage{20}
%<<CONT
divisor of $P$ is presupposed, and from the representation (10)
of the group $M$ it follows that $M = Q/Q_1$.

It remains to show that if $M$ is a normal divisor of
$P/Q_1$, then $Q$ is also a normal divisor of $P$. This follows thus:

If $e$ is any element in $P$ and $Q_1 b_k$ an element in $M$,
then from the assumption that $M$ is a normal divisor of $P/Q_1$ it follows that:

\[ Q_1 e^{-1} Q_1 b_k Q_1 e = Q_1 b_h, \]

which we can also express, according to the composition of the parts, by the formula

\weq{14}{Q_1 e^{-1} Q_1 B Q_1 e = Q_1 B}

Now since $Q_1$ is a normal divisor of $P$,
$Q_1$ can be interchanged with each of the other factors, so
that from (14) follows:

\[ e^{-1} Q_1 B e = Q_1 B, \]

or

\weq{15}{e^{-1} Q e = Q,}

whereby it is expressed that $Q$ is a normal divisor of $P$.

From Theorems 1. and 2. results the corollary:

\gesperrt{3. A normal divisor $Q$ of $P$ is a greatest
normal divisor if and only if the complementary
 group $P/Q$ is simple.}

If $Q$ and $Q'$ are two normal divisors of $P$, then the
symbolic product $Q Q'$ is also a normal divisor.

For, since $Q$ is a normal divisor, we have

\[ Q Q' = Q' Q, \]

hence

\[ Q Q' Q Q' = Q Q Q' Q' = Q Q', \]

and this, by \S.\,4, 1., is the criterion that $Q Q'$ is a group.
That it is a normal divisor of $P$ then follows
likewise, by \S.\,4, 2., from

\[ Q Q' B = Q B Q' = B Q Q', \]

where $B$ is any part of $P$.

Now if $Q$ is a \gesperrt{greatest normal divisor} of $P$, and $Q'$ is
different from $Q$, then $Q Q'$ is a normal divisor of $P$ which
contains both $Q$ and $Q'$ within itself and hence is different from $Q$.
Consequently, since above $Q$ there is no normal divisor of $P$
except $P$ itself,

\weq{16}{Q Q' = P,}

and likewise we must also have

\weq{17}{Q' Q = P}
%CONT>>
